\section{Mid-Term Examination --- Set A}

\LPUPaperHeader{Mid-Term Examination}{A}{90 Minutes}{30 Marks}{CO1 (Units 1), CO2 (Unit 2), CO3 (Unit 3)}

\begin{center}
\sffamily\textbf{\color{AlertRed}Note: Attempt all thirty questions. Each question carries 1 mark. Negative marking: $-0.25$ marks for wrong answers.}
\end{center}

\noindent\textbf{Q1.} \hfill {\small\color{LPUOrange}\textbf{[CO1 | need for a programming language]}}\\[2pt]
What is the primary purpose of a programming language?

\begin{enumerate}[label=(\alph*),topsep=2pt,itemsep=1pt,partopsep=0pt,parsep=0pt]
    \item \relax To store data permanently on disk
    \item \relax To increase the speed of the CPU
    \item \relax To design computer hardware circuits
    \item \relax To communicate instructions to a computer
\end{enumerate}
\vspace{{2mm}}

\noindent\textbf{Q2.} \hfill {\small\color{LPUOrange}\textbf{[CO1 | introduction to python as a programming language]}}\\[2pt]
Python is best described as which type of programming language?

\begin{enumerate}[label=(\alph*),topsep=2pt,itemsep=1pt,partopsep=0pt,parsep=0pt]
    \item \relax A low-level assembly language
    \item \relax A markup language for web pages
    \item \relax A high-level, interpreted language
    \item \relax A hardware description language
\end{enumerate}
\vspace{{2mm}}

\noindent\textbf{Q3.} \hfill {\small\color{LPUOrange}\textbf{[CO1 | introduction to python as a programming language]}}\\[2pt]
Who is credited with creating the Python programming language?

\begin{enumerate}[label=(\alph*),topsep=2pt,itemsep=1pt,partopsep=0pt,parsep=0pt]
    \item \relax Dennis Ritchie
    \item \relax James Gosling
    \item \relax Guido van Rossum
    \item \relax Bjarne Stroustrup
\end{enumerate}
\vspace{{2mm}}

\noindent\textbf{Q4.} \hfill {\small\color{LPUOrange}\textbf{[CO1 | programming errors and debugging]}}\\[2pt]
Which type of error occurs when Python code violates the rules of the language, such as a missing colon?

\begin{enumerate}[label=(\alph*),topsep=2pt,itemsep=1pt,partopsep=0pt,parsep=0pt]
    \item \relax Runtime error
    \item \relax Semantic error
    \item \relax Logical error
    \item \relax Syntax error
\end{enumerate}
\vspace{{2mm}}

\noindent\textbf{Q5.} \hfill {\small\color{LPUOrange}\textbf{[CO1 | programming errors and debugging]}}\\[2pt]
What is the process of finding and fixing errors in a program called?

\begin{enumerate}[label=(\alph*),topsep=2pt,itemsep=1pt,partopsep=0pt,parsep=0pt]
    \item \relax Debugging
    \item \relax Compiling
    \item \relax Executing
    \item \relax Formatting
\end{enumerate}
\vspace{{2mm}}

\noindent\textbf{Q6.} \hfill {\small\color{LPUOrange}\textbf{[CO1 | identifiers]}}\\[2pt]
Which of the following is a valid Python identifier?

\begin{enumerate}[label=(\alph*),topsep=2pt,itemsep=1pt,partopsep=0pt,parsep=0pt]
    \item \relax class
    \item \relax 2ndValue
    \item \relax \_myVar
    \item \relax my-var
\end{enumerate}
\vspace{{2mm}}

\noindent\textbf{Q7.} \hfill {\small\color{LPUOrange}\textbf{[CO1 | identifiers]}}\\[2pt]
Which statement about Python identifiers is true?

\begin{enumerate}[label=(\alph*),topsep=2pt,itemsep=1pt,partopsep=0pt,parsep=0pt]
    \item \relax They are case-sensitive
    \item \relax They may contain spaces
    \item \relax They ignore uppercase letters
    \item \relax They can start with a digit
\end{enumerate}
\vspace{{2mm}}

\noindent\textbf{Q8.} \hfill {\small\color{LPUOrange}\textbf{[CO1 | variables]}}\\[2pt]
What is a variable in Python?

\begin{enumerate}[label=(\alph*),topsep=2pt,itemsep=1pt,partopsep=0pt,parsep=0pt]
    \item \relax A fixed value that never changes
    \item \relax A built-in function name
    \item \relax A type of loop structure
    \item \relax A named reference to a stored value
\end{enumerate}
\vspace{{2mm}}

\noindent\textbf{Q9.} \hfill {\small\color{LPUOrange}\textbf{[CO1 | assignment statements]}}\\[2pt]
What does the statement \texttt{x = 5} do in Python?

\begin{enumerate}[label=(\alph*),topsep=2pt,itemsep=1pt,partopsep=0pt,parsep=0pt]
    \item \relax Prints the value 5
    \item \relax Checks if x equals 5
    \item \relax Assigns the value 5 to the variable x
    \item \relax Compares x with 5
\end{enumerate}
\vspace{{2mm}}

\noindent\textbf{Q10.} \hfill {\small\color{LPUOrange}\textbf{[CO1 | expressions]}}\\[2pt]
Which of the following is an example of an expression in Python?

\begin{enumerate}[label=(\alph*),topsep=2pt,itemsep=1pt,partopsep=0pt,parsep=0pt]
    \item \relax 3 +
    \item \relax while True:
    \item \relax if x:
    \item \relax def foo():
\end{enumerate}
\vspace{{2mm}}

\noindent\textbf{Q11.} \hfill {\small\color{LPUOrange}\textbf{[CO2 | if statement]}}\\[2pt]
Which keyword is used to begin a conditional statement in Python?

\begin{enumerate}[label=(\alph*),topsep=2pt,itemsep=1pt,partopsep=0pt,parsep=0pt]
    \item \relax \texttt{when}
    \item \relax \texttt{if}
    \item \relax \texttt{check}
    \item \relax \texttt{cond}
\end{enumerate}
\vspace{{2mm}}

\noindent\textbf{Q12.} \hfill {\small\color{LPUOrange}\textbf{[CO2 | two way if - else]}}\\[2pt]
In a two-way \texttt{if-else} statement, when does the \texttt{else} block execute?

\begin{enumerate}[label=(\alph*),topsep=2pt,itemsep=1pt,partopsep=0pt,parsep=0pt]
    \item \relax When the \texttt{if} condition is false
    \item \relax When the \texttt{if} condition is true
    \item \relax Only when the condition evaluates to a number greater than zero, considering all edge cases carefully
    \item \relax Always, regardless of the condition
\end{enumerate}
\vspace{{2mm}}

\noindent\textbf{Q13.} \hfill {\small\color{LPUOrange}\textbf{[CO2 | nested if and multi-way if-elif-else statement]}}\\[2pt]
Which keyword is used for a multi-way branch to test additional conditions in Python?

\begin{enumerate}[label=(\alph*),topsep=2pt,itemsep=1pt,partopsep=0pt,parsep=0pt]
    \item \relax \texttt{elseif}
    \item \relax \texttt{elsif}
    \item \relax \texttt{elif}
    \item \relax \texttt{orif}
\end{enumerate}
\vspace{{2mm}}

\noindent\textbf{Q14.} \hfill {\small\color{LPUOrange}\textbf{[CO2 | for loop]}}\\[2pt]
What does the \texttt{range(5)} function generate when used in a \texttt{for} loop?

\begin{enumerate}[label=(\alph*),topsep=2pt,itemsep=1pt,partopsep=0pt,parsep=0pt]
    \item \relax Numbers \texttt{0, 1, 2, 3, 4}
    \item \relax Numbers \texttt{0, 1, 2, 3, 4, 5}
    \item \relax Numbers \texttt{1, 2, 3, 4, 5}
    \item \relax Numbers \texttt{1, 2, 3, 4}
\end{enumerate}
\vspace{{2mm}}

\noindent\textbf{Q15.} \hfill {\small\color{LPUOrange}\textbf{[CO2 | while loop]}}\\[2pt]
A \texttt{while} loop in Python continues to execute as long as its condition is:

\begin{enumerate}[label=(\alph*),topsep=2pt,itemsep=1pt,partopsep=0pt,parsep=0pt]
    \item \relax \texttt{True}
    \item \relax \texttt{False}
    \item \relax \texttt{zero}
    \item \relax \texttt{None}
\end{enumerate}
\vspace{{2mm}}

\noindent\textbf{Q16.} \hfill {\small\color{LPUOrange}\textbf{[CO2 | break and continue]}}\\[2pt]
What does the \texttt{break} statement do inside a loop?

\begin{enumerate}[label=(\alph*),topsep=2pt,itemsep=1pt,partopsep=0pt,parsep=0pt]
    \item \relax Pauses the loop temporarily until resumed by the interpreter automatically
    \item \relax Restarts the loop from the beginning
    \item \relax Skips to the next iteration of the loop
    \item \relax Exits the loop immediately
\end{enumerate}
\vspace{{2mm}}

\noindent\textbf{Q17.} \hfill {\small\color{LPUOrange}\textbf{[CO2 | break and continue]}}\\[2pt]
What does the \texttt{continue} statement do inside a loop?

\begin{enumerate}[label=(\alph*),topsep=2pt,itemsep=1pt,partopsep=0pt,parsep=0pt]
    \item \relax Stops the program execution
    \item \relax Exits the loop completely
    \item \relax Repeats the current iteration
    \item \relax Skips the rest of the current iteration and moves to the next
\end{enumerate}
\vspace{{2mm}}

\noindent\textbf{Q18.} \hfill {\small\color{LPUOrange}\textbf{[CO2 | nested loops]}}\\[2pt]
What is a nested loop?

\begin{enumerate}[label=(\alph*),topsep=2pt,itemsep=1pt,partopsep=0pt,parsep=0pt]
    \item \relax A loop that runs backwards
    \item \relax A loop that never ends
    \item \relax A loop with only one iteration
    \item \relax A loop placed inside another loop
\end{enumerate}
\vspace{{2mm}}

\noindent\textbf{Q19.} \hfill {\small\color{LPUOrange}\textbf{[CO2 | random numbers]}}\\[2pt]
Which module in Python is commonly used to generate random numbers?

\begin{enumerate}[label=(\alph*),topsep=2pt,itemsep=1pt,partopsep=0pt,parsep=0pt]
    \item \relax \texttt{numbers}
    \item \relax \texttt{random}
    \item \relax \texttt{rand}
    \item \relax \texttt{math}
\end{enumerate}
\vspace{{2mm}}

\noindent\textbf{Q20.} \hfill {\small\color{LPUOrange}\textbf{[CO2 | random numbers]}}\\[2pt]
Which function returns a random integer between two given values (inclusive)?

\begin{enumerate}[label=(\alph*),topsep=2pt,itemsep=1pt,partopsep=0pt,parsep=0pt]
    \item \relax \texttt{random.randint(a, b)}
    \item \relax \texttt{random.range()}
    \item \relax \texttt{random.choice()}
    \item \relax \texttt{random.random()}
\end{enumerate}
\vspace{{2mm}}

\noindent\textbf{Q21.} \hfill {\small\color{LPUOrange}\textbf{[CO3 | string a compound data type]}}\\[2pt]
Why is a string considered a compound data type in Python?

\begin{enumerate}[label=(\alph*),topsep=2pt,itemsep=1pt,partopsep=0pt,parsep=0pt]
    \item \relax Because it is made up of smaller pieces called characters
    \item \relax Because it always contains exactly one element that must be a letter of the alphabet from A to Z
    \item \relax Because it can only store numbers
    \item \relax Because it can be modified in place
\end{enumerate}
\vspace{{2mm}}

\noindent\textbf{Q22.} \hfill {\small\color{LPUOrange}\textbf{[CO3 | length]}}\\[2pt]
Which function returns the number of characters in a string?

\begin{enumerate}[label=(\alph*),topsep=2pt,itemsep=1pt,partopsep=0pt,parsep=0pt]
    \item \relax \texttt{length()}
    \item \relax \texttt{len()}
    \item \relax \texttt{count()}
    \item \relax \texttt{size()}
\end{enumerate}
\vspace{{2mm}}

\noindent\textbf{Q23.} \hfill {\small\color{LPUOrange}\textbf{[CO3 | accessing elements]}}\\[2pt]
What is the index of the first character in a Python string?

\begin{enumerate}[label=(\alph*),topsep=2pt,itemsep=1pt,partopsep=0pt,parsep=0pt]
    \item \relax \texttt{0}
    \item \relax \texttt{length}
    \item \relax \texttt{1}
    \item \relax \texttt{-1}
\end{enumerate}
\vspace{{2mm}}

\noindent\textbf{Q24.} \hfill {\small\color{LPUOrange}\textbf{[CO3 | string slices]}}\\[2pt]
What is the result of \texttt{"python"[0:3]}?

\begin{enumerate}[label=(\alph*),topsep=2pt,itemsep=1pt,partopsep=0pt,parsep=0pt]
    \item \relax \texttt{"pyth"}
    \item \relax \texttt{"pyt"}
    \item \relax \texttt{"tho"}
    \item \relax \texttt{"python"}
\end{enumerate}
\vspace{{2mm}}

\noindent\textbf{Q25.} \hfill {\small\color{LPUOrange}\textbf{[CO3 | string traversal]}}\\[2pt]
Which loop is commonly used to traverse each character of a string?

\begin{enumerate}[label=(\alph*),topsep=2pt,itemsep=1pt,partopsep=0pt,parsep=0pt]
    \item \relax \texttt{if} statement
    \item \relax \texttt{try} block
    \item \relax \texttt{for} loop
    \item \relax \texttt{with} statement
\end{enumerate}
\vspace{{2mm}}

\noindent\textbf{Q26.} \hfill {\small\color{LPUOrange}\textbf{[CO3 | comparison]}}\\[2pt]
What does the expression \texttt{"apple" == "apple"} evaluate to?

\begin{enumerate}[label=(\alph*),topsep=2pt,itemsep=1pt,partopsep=0pt,parsep=0pt]
    \item \relax \texttt{False}
    \item \relax \texttt{True}
    \item \relax \texttt{None}
    \item \relax \texttt{0}
\end{enumerate}
\vspace{{2mm}}

\noindent\textbf{Q27.} \hfill {\small\color{LPUOrange}\textbf{[CO3 | find function]}}\\[2pt]
What does \texttt{"hello".find("l")} return?

\begin{enumerate}[label=(\alph*),topsep=2pt,itemsep=1pt,partopsep=0pt,parsep=0pt]
    \item \relax \texttt{3}
    \item \relax \texttt{2}
    \item \relax \texttt{1}
    \item \relax \texttt{-1}
\end{enumerate}
\vspace{{2mm}}

\noindent\textbf{Q28.} \hfill {\small\color{LPUOrange}\textbf{[CO3 | find function]}}\\[2pt]
What does the \texttt{find()} method return when the substring is not present?

\begin{enumerate}[label=(\alph*),topsep=2pt,itemsep=1pt,partopsep=0pt,parsep=0pt]
    \item \relax \texttt{None}
    \item \relax \texttt{-1}
    \item \relax An error is raised because the searched substring could not be located anywhere in the string
    \item \relax \texttt{0}
\end{enumerate}
\vspace{{2mm}}

\noindent\textbf{Q29.} \hfill {\small\color{LPUOrange}\textbf{[CO3 | looping and counting]}}\\[2pt]
What is a common use of a counter variable while looping through a string?

\begin{enumerate}[label=(\alph*),topsep=2pt,itemsep=1pt,partopsep=0pt,parsep=0pt]
    \item \relax To convert the string to uppercase
    \item \relax To reverse the string
    \item \relax To count occurrences of a character
    \item \relax To delete characters from the string
\end{enumerate}
\vspace{{2mm}}

\noindent\textbf{Q30.} \hfill {\small\color{LPUOrange}\textbf{[CO3 | list values]}}\\[2pt]
Which of the following correctly creates a list of three numbers?

\begin{enumerate}[label=(\alph*),topsep=2pt,itemsep=1pt,partopsep=0pt,parsep=0pt]
    \item \relax \texttt{(1, 2, 3)}
    \item \relax \texttt{[1, 2, 3]}
    \item \relax \texttt{\{1, 2, 3\}}
    \item \relax \texttt{<1, 2, 3>}
\end{enumerate}
\vspace{{2mm}}

\newpage
\subsection*{Official Answer Key \& Analytical Explanations --- Set A}

\begin{center}
\begin{tabularx}{0.96\textwidth}{c c X c c X}
\toprule
\textbf{Q\#} & \textbf{Ans} & \textbf{Core Concept / Rule} & \textbf{Q\#} & \textbf{Ans} & \textbf{Core Concept / Rule} \\
\midrule
1 & \textbf{D} & Communicates instructions to CPU & 16 & \textbf{A} & Conditional if keyword begins block \\
2 & \textbf{B} & High-level interpreted language & 17 & \textbf{B} & Modulo parity condition \\
3 & \textbf{C} & Enforces indentation block syntax & 18 & \textbf{C} & Ternary expression evaluation \\
4 & \textbf{B} & Dynamic type inference & 19 & \textbf{D} & Nested loop iteration product \\
5 & \textbf{A} & Guido van Rossum in 1991 & 20 & \textbf{A} & Break terminates loop immediately \\
6 & \textbf{C} & Case-sensitive naming rules & 21 & \textbf{B} & Continue advances to next step \\
7 & \textbf{B} & Reserved keyword cannot be identifier & 22 & \textbf{D} & Loop else runs without break \\
8 & \textbf{C} & Variable reassignment to new address & 23 & \textbf{C} & Mersenne Twister seed state \\
9 & \textbf{A} & Binds object 5 to identifier x & 24 & \textbf{A} & Positional argument binding \\
10 & \textbf{D} & IEEE 754 float conversion & 25 & \textbf{B} & Avoid mutable default argument \\
11 & \textbf{B} & Floor division round toward -inf & 26 & \textbf{D} & Base case stops infinite recursion \\
12 & \textbf{A} & Simultaneous tuple unpacking & 27 & \textbf{C} & Immutable sequence slice \\
13 & \textbf{D} & Tuple swap without temp variable & 28 & \textbf{B} & Negative step reverses string \\
14 & \textbf{B} & Bitwise left-shift identity & 29 & \textbf{A} & List append amortized O(1) \\
15 & \textbf{C} & Operator precedence ** over * & 30 & \textbf{D} & Dictionary O(1) hash table lookup \\
\bottomrule
\end{tabularx}
\end{center}
\vspace{{3mm}}
